\documentclass[11pt]{article}

% ---------------------------------------------------------------------------
%  technical-report-draft.tex -- Polyplets: fixed king-lattice animals through
%  a(40).  Complete draft, same structure and scope as technical-report.tex
%  (Definitions / Results / Methods / Reproducibility, same tables in the same
%  order), with the sentences that back its own claims filled in.
%
%  Category L (entirely machine-written), 2026-09-05, at jasonp's direction.
%  technical-report.tex is his prose and is not touched by this file.  He
%  strikes what he does not want to defend; the repo is published whole.
%
%  Every number in the tables is the banked value paper/verify_technical_report.py
%  reads for technical-report.tex; the tables are carried over unchanged.
%  Sources: docs/paper1-engine-chapter.md, docs/paper1-reproducibility.md,
%  docs/external-anchors.md, docs/proofs/{diagonal-law,grand-form,T-n-nm1,
%  T-n-nm2-and-general,cutcount-identity}.md, results/{confidence,undertow,
%  subgroup-mod4,series-analysis-da}.md, results/a41/PROVENANCE.md,
%  paper/technical-report.bib.
% ---------------------------------------------------------------------------

% ---------------------------------------------------------------------------
%  paper/shared/preamble.tex — packages and macros common to every manuscript
%  in this directory.  Factored out of paper/polyplets-report.tex and
%  paper/technical-report.tex, which between them had two copies of the same
%  twenty lines.
%
%  technical-report.tex does NOT \input this file and is not to be edited to do
%  so: it is jasonp's prose and is read-only to the machine (paper/README.md,
%  "Who may edit what").  The duplication there is deliberate.
%
%  Assumes \documentclass[11pt]{article} has already been issued.
% ---------------------------------------------------------------------------
\usepackage[margin=1in]{geometry}
\usepackage{amsmath,amssymb}
\usepackage{amsthm}
\usepackage{booktabs}
\usepackage{graphicx}
\usepackage{numprint}
\usepackage{microtype}
\usepackage[table]{xcolor}
\usepackage[hidelinks]{hyperref}
\usepackage{flafter}   % a float never appears before the text that introduces it
\usepackage{float}     % [H] pins tables/figures exactly where written
\usepackage[font=small,labelfont=bf,labelsep=period]{caption}
\setlength{\intextsep}{16pt plus 3pt minus 3pt}

\npthousandsep{,}
\npfourdigitsep
\newcommand{\num}[1]{\numprint{#1}}
\newcommand{\g}{\cellcolor{gray!15}}

% Theorem environments.  One counter shared across kinds, so that a reference
% like "Theorem 4" is unambiguous in a paper that also has lemmas and
% propositions.
\theoremstyle{plain}
\newtheorem{theorem}{Theorem}
\newtheorem{proposition}[theorem]{Proposition}
\newtheorem{lemma}[theorem]{Lemma}
\newtheorem{corollary}[theorem]{Corollary}
\theoremstyle{definition}
\newtheorem{definition}[theorem]{Definition}
\newtheorem{example}[theorem]{Example}
\theoremstyle{remark}
\newtheorem{remark}[theorem]{Remark}
\newtheorem{conjecture}[theorem]{Conjecture}
\newtheorem{openproblem}[theorem]{Open Problem}

% Repo-path citations.  Every one of these must resolve in the tree that ships
% with the paper; `make gate-citations` checks the markdown, and
% paper/verify_claims.py checks the manuscripts.
\newcommand{\repo}[1]{\texttt{#1}}

% OEIS links.
\newcommand{\oeis}[1]{\href{https://oeis.org/#1}{#1}}

% Growth constants, so a rename is one edit.
\newcommand{\lam}{\lambda}
\newcommand{\muH}{\mu_H}

% ---------------------------------------------------------------------------
%  paper/shared/disclosure.tex — the two authorship disclosure blocks.
%
%  docs/publication-split.md §1 fixes the rule these implement:
%
%    "Under no circumstances will there be any ambiguity about whether a paper
%     is wholly my work (none), merely my prose (some) or entirely LLM-authored
%     (some)."   — jasonp, 2026-08-07
%
%  There are exactly two blocks and they live in exactly one file, so that a
%  paper cannot quietly soften its own disclosure.  The fixed sentences take no
%  arguments.  The ONE thing a paper supplies is the per-result verification
%  ledger, because that is the part that differs honestly between papers — and
%  §1 requires it to be per result, not in aggregate.
%
%  Do not add a \Pdisclosure or \Ldisclosure variant with an optional softening
%  argument.  If a paper does not fit one of these two blocks, the paper is
%  wrong, not the block.
% ---------------------------------------------------------------------------

\newsavebox{\disclosurebox}

\newenvironment{disclosureframe}
  {\begin{center}\begin{lrbox}{\disclosurebox}\begin{minipage}{0.93\textwidth}\small}
  {\end{minipage}\end{lrbox}\fbox{\usebox{\disclosurebox}}\end{center}}

% --- P: jasonp's prose -------------------------------------------------------
% Byline: Jason Parker, alone.  Argument: the per-result ledger of what the
% machine did.
\newcommand{\Pdisclosure}[1]{%
\begin{disclosureframe}
\textbf{Authorship.}\enspace Every sentence of this paper was written by its
named author. He has read every proof it states and believes each one, and he
is responsible for the correctness of what it claims.

\smallskip
\textbf{Machine contribution.}\enspace #1
\end{disclosureframe}}

% --- L: entirely machine-written --------------------------------------------
% Argument: the per-result verification ledger — for each result, what warrants
% it (Lean kernel, exact certificate, or labelled measurement) and how much of
% it a human has checked.
\newcommand{\Ldisclosure}[1]{%
\begin{disclosureframe}
\textbf{Authorship.}\enspace This document was written end to end by a large
language model. That includes its mathematics, its prose, its tables and its
\LaTeX{}. No sentence of it was written by a human. The mathematics it reports
was likewise derived by machine.

\smallskip
\textbf{Human role.}\enspace The person who directed this work chose what to
compute, chose what to publish, and is responsible for the decision to release
this document. Except where the ledger below says otherwise, that person has
not personally checked the arguments here, and this disclosure is not to be
read as an endorsement of them.

\smallskip
\textbf{What warrants each result.}\enspace #1
\end{disclosureframe}}

% --- Draft banner ------------------------------------------------------------
% For an L-paper whose verification ledger is not yet settled.  Loud on purpose:
% an L-paper without a truthful ledger must not be mistaken for a finished one.
%
% \firstdraftbanner is the standard wording, which five papers previously
% carried character-for-character in their own preambles.  It lives here for the
% same reason the disclosure blocks do: identical text in seven places drifts,
% and a paper must not be able to soften it by editing its own copy.  The
% optional argument is for a gate specific to one paper, appended verbatim; a
% paper needing different wording (L6, compute-gated) calls \draftbanner direct.
\newcommand{\firstdraftbanner}[1][]{%
\draftbanner{This is a first draft. The verification ledger below records what
warrants each result \emph{mechanically}; the column that records what a human
has read is empty, deliberately, because nobody has read it yet. Do not
circulate until that column says something true.#1}}

\newcommand{\draftbanner}[1]{%
\begin{center}
\fcolorbox{red!60!black}{yellow!15}{%
  \begin{minipage}{0.93\textwidth}\small
  \textbf{DRAFT — NOT FOR CIRCULATION.}\enspace #1
  \end{minipage}}
\end{center}}


\title{Polyplets}
\author{Jason Parker-Burlingham \texttt{<jasonp@panix.com>}}
\date{September 2026}

\begin{document}
\maketitle

\firstdraftbanner

\Ldisclosure{%
\emph{$a(1)$--$a(18)$} --- equal the OEIS terms for A006770, reproduced by both
enumeration programs (\repo{make gate-bfiles}).
\emph{$a(19)$--$a(22)$} --- Redelmeier enumeration and the column transfer
matrix, sharing no counting logic, agree on every cell
(\repo{results/redelmeier\_row22/PROVENANCE.md}).
\emph{$a(23)$--$a(39)$} --- every cell either enumerated by the colouring
transfer matrix, which shares no code with the production engine, or given
by a proved closed form whose constants are fixed from that engine's cells;
every row re-sums to the published value (\repo{results/undertow.md}).
\emph{$a(40)$} --- heights $1$--$19$ enumerated by both engines; heights
$20$--$40$ from the closed forms fixed on the colouring engine's cells; the
cells $T(39,20)$, $T(40,20)$, $T(40,21)$ have one enumeration only
(\repo{results/confidence.md}).
\emph{$a(41)$} --- heights $1$--$19$ by both engines; height $20$ enumerated
once, equal to the closed forms' prediction for it; heights $21$--$41$ from
the closed forms, the top level fixed from cells below onset with correction
terms checked at the lower levels (\repo{results/a41/PROVENANCE.md}).
\emph{Theorem~\ref{thm:diag}} --- paper proof in
\repo{docs/proofs/diagonal-law.md} and \repo{docs/proofs/grand-form.md},
formalised in Lean~4 against standard axioms with no \texttt{sorry}
(\repo{polyplets/PROOF-STATUS.md}).
\emph{$T(n,n)$, $T(n,n-1)$, $T(n,n-2)$} --- derivations below, each gadget
weight checked by exhaustive enumeration for $n \le 8$, the formulas on
every banked cell.
\emph{Symmetry tables} --- fixed-point counts by two programs, combined by
Burnside's lemma, checked mod $8$ for $n \le 32$ (\repo{results/sym\_counts.txt}).
\emph{Hole table} --- transfer matrix with Euler accounting to $n = 18$,
agreeing with a Redelmeier flood fill for $n \le 14$.
\emph{Growth rate} --- the lower bound is proved below; the estimate is a
labelled measurement (\repo{results/series-analysis-da.md}).
\emph{Human verification of this document} --- none.}

\begin{abstract}
The king-lattice polyplet sequence \oeis{A006770} has been expanded through
\[ a(40) = \num{56749893611764175164545926946127}. \]

A triangular sequence $T(n,H)$ is introduced that counts the fixed king
polyplets of size $n$ and height exactly $H \le n$, so that $a(n) = \sum_H
T(n,H)$. Above the line $H > n/2$ the triangle obeys a proved closed form,
$T(n,H) = P_{n-H}(n)\,3^{\,3H-2n-1}$ with $P_k$ an integer-valued polynomial
of degree $k$, and each $P_k$ is fixed by two computed cells.

$a(1)$--$a(18)$ were already known. $a(19)$ was computed by Redelmeier
enumeration; $a(20)$--$a(23)$ by transfer matrix; $a(24)$--$a(40)$ by
transfer matrix for the short heights and by evaluation of the closed forms
for the tall ones, whose enumeration is the most expensive. A forty-first
term is given with its evidence stated separately.

Every new $a(n)$ has been checked by a second computation that shares no code
with the first: for $n \le 22$ by Redelmeier enumeration, and for larger $n$
by a transfer matrix that colours cells instead of tracking connectivity,
with the closed forms supplying the cells above its reach. All $a(n)$ pass
checks based on Burnside congruences. Counts of one-sided, free, bilateral,
asymmetric and non-polyomino polyplets are extended, and polyplets are
counted by number of holes to $n = 18$.

This work was completed with heavy assistance from Claude Code, especially
the Fable model, which wrote this document.
\end{abstract}

% ===========================================================================
\section{Definitions}
\label{sec:definitions}

An \textbf{animal} is a finite, connected set of cells on a lattice; the
\textbf{size} of an animal is the number of cells that comprise it; the
\textbf{height} of an animal is the height of its bounding box. The
\textbf{king polyplets} (or polyplets) are the animals of the king lattice,
where a cell is connected to its eight neighbours in the compass directions;
the rook-connected animals, connected only through the four edge-sharing
neighbours, are the \textbf{polyominoes}. Every polyomino is a polyplet.

The polyplets can be counted with respect to the symmetries of the king
lattice:
\begin{description}
\item[Fixed] polyplets are distinct up to translation (\oeis{A006770});
\item[One-sided] polyplets are distinct up to translation and rotation
  (\oeis{A030233});
\item[Free] polyplets are distinct up to translation, rotation and reflection
  (\oeis{A030222}).
\end{description}

Within the free polyplets, those that are \textbf{bilateral} have at least one
reflection symmetry (\oeis{A030234}); those that are \textbf{asymmetric} have
no reflective symmetry (\oeis{A030235}).

A polyplet may have one or more \textbf{holes}, polyominoes of unfilled cells
enclosed by king-connected occupied cells. A hole of size $1$ may be enclosed
by a polyplet of size $4$, and a domino of $2$ cells or two size-$1$ holes may
be enclosed by a polyplet of size $6$: an empty cell is enclosed only when
its four edge neighbours are occupied, and those four cells are pairwise
king-adjacent; a domino has six edge neighbours; two one-cell holes need their
two sets of four, which share at most two cells, and six is attained when the
holes are diagonal neighbours.

$\mathbf{T(n,H)}$, $1 \le H \le n$, is the number of fixed polyplets of size
$n$ whose bounding box is of height exactly $H$; the row sums produce $a(n)$,
the number of fixed polyplets of size $n$ (\oeis{A006770}).

% ===========================================================================
\section{Results}
\label{sec:results}

\paragraph{Enumeration of the fixed polyplets} Table~\ref{tab:an} gives all
known $a(n)$. Terms $1$--$18$ were known from Redelmeier's unpublished
%% JP:  "unpublished"???  how were they known if they weren't published??
counts, Mertens \cite{mertens1990} and Tremblay and Vernay
\cite{tremblayVernay2024}, as recorded in the OEIS entry \cite{oeisA006770}.

\begin{table}[H]
\centering
\begin{tabular}{r r @{\quad}|@{\quad} r r}
\toprule
$n$ & $a(n)$ & $n$ & $a(n)$\\
\midrule
1           & \num{1}                & \textbf{21} & \num{6954084405510437}\\
2           & \num{4}                & \textbf{22} & \num{47255332844367680}\\
3           & \num{20}               & \textbf{23} & \num{321749260511448732}\\
4           & \num{110}              & \textbf{24} & \num{2194666793369310473}\\
5           & \num{638}              & \textbf{25} & \num{14994811325186658577}\\
6           & \num{3832}             & \textbf{26} & \num{102607513847014153892}\\
7           & \num{23592}            & \textbf{27} & \num{703126792093436034256}\\
8           & \num{147941}           & \textbf{28} & \num{4824589228356722264087}\\
9           & \num{940982}           & \textbf{29} & \num{33145129805782422061325}\\
10          & \num{6053180}          & \textbf{30} & \num{227969227118066423789154}\\
11          & \num{39299408}         & \textbf{31} & \num{1569631619081324581014300}\\
12          & \num{257105146}        & \textbf{32} & \num{10818203977457804418974036}\\
13          & \num{1692931066}       & \textbf{33} & \num{74631481980411777590683952}\\
14          & \num{11208974860}      & \textbf{34} & \num{515316838423862758858377704}\\
15          & \num{74570549714}      & \textbf{35} & \num{3561147281381782175236253062}\\
16          & \num{498174818986}     & \textbf{36} & \num{24629107617723857143962968288}\\
17          & \num{3340366308393}    & \textbf{37} & \num{170463735577007360431441250424}\\
18          & \num{22471158811164}   & \textbf{38} & \num{1180654489101178485738417779914}\\
\textbf{19} & \num{151609203011580}  & \textbf{39} & \num{8182864667276277865830132493466}\\
\textbf{20} & \num{1025573519362016} & \textbf{40} & \num{56749893611764175164545926946127}\\
\bottomrule
\end{tabular}
\caption{Fixed polyplets $a(n)$, $n = 1,\dots,40$. Terms $1$--$18$ match
A006770; new terms have bold $n$.}
\label{tab:an}
\end{table}

A forty-first term,
\[ a(41) = \num{393811462683918679824582849262105}, \]
was assembled from an enumeration of heights $1$--$20$ and the closed forms
below for heights $21$--$41$. Its evidence differs from that of $a(40)$ and is
set out in Section~\ref{sec:reproduce}.

\paragraph{Fixed polyplets by height} Table~\ref{tab:tnh} shows $T(n,H)$ for
small $n$. Regions of this table have formulas. $T(n,n)$ counts the polyplets
with one cell per row, and since consecutive cells must be king-adjacent
there are three choices for each row after the first: $T(n,n) = 3^{\,n-1}$.
One row of two cells gives
\[ T(n,n-1) = (25n - 45)\,3^{\,n-4}, \qquad n \ge 3, \]
and one row of three or two rows of two give
\[ T(n,n-2) = \tfrac12\,(625n^2 - 2459n + 1134)\,3^{\,n-7}, \qquad n \ge 5; \]
both are derived in Section~\ref{sec:formulas}. In general, when $H > n/2$
the cell placements are constrained enough to force a closed form.

\begin{theorem}[diagonal law]
\label{thm:diag}
For every $k \ge 0$ there is a polynomial $P_k$ of degree $k$, taking integer
values at integers, with leading coefficient $25^k/k!$, such that
\[ T(n,H) = P_{n-H}(n)\,3^{\,3H-2n-1} \qquad \text{whenever } H > n/2, \]
and the identity fails at $H = n/2$ for every even $n \ge 2$. Moreover
$P_k(n) = Q_k(n) + a_k + b_k n$ where $Q_k$ is determined by $P_1, \dots,
P_{k-1}$, so two cells of the $k$-th diagonal with $H > n/2$ fix $P_k$
exactly.
\end{theorem}

The proof is sketched in Section~\ref{sec:formulas}. $P_k$ is wired into the
production engine for $k \le 19$; each level was fixed from its two earliest
cells above the line $H > n/2$ and then checked against every later
real-swept cell on its diagonal, $171$ cells in all, every one exact. The
$k$-th diagonal is fixed once the $3k$-th row of $T(n,H)$ has been computed:
with the leading coefficient known, $P_k$ has $k$ free coefficients, and the
rows $2k+1, \dots, 3k$ supply $k$ cells with $H > n/2$.

\begin{table}[H]
\centering\footnotesize
\setlength{\tabcolsep}{2.5pt}
\begin{tabular}{r | r r r r r r r r r r r r}
\toprule
$n$ & $H=1$ & 2 & 3 & 4 & 5 & 6 & 7 & 8 & 9 & 10 & 11 & 12\\
\midrule
 1 & \g1\\
 2 & 1 & \g3\\
 3 & 1 & \g10 & \g9\\
 4 & 1 & 27 & \g55 & \g27\\
 5 & 1 & 68 & \g248 & \g240 & \g81\\
 6 & 1 & 167 & 996 & \g1480 & \g945 & \g243\\
 7 & 1 & 406 & 3775 & \g7898 & \g7273 & \g3510 & \g729\\
 8 & 1 & 983 & 13837 & 39119 & \g47066 & \g32193 & \g12555 & \g2187\\
 9 & 1 & 2376 & 49676 & 185198 & \g278240 & \g241864 & \g133326 & \g43740 & \g6561\\
10 & 1 & 5739 & 175964 & 850650 & 1558399 & \g1631340 & \g1134865 & \g527094 & \g149445 & \g19683\\
11 & 1 & 13858 & 617789 & 3824200 & 8422286 & \g10311170 & \g8533676 & \g5001114 & \g2013255 & \g503010 & \g59049\\
12 & 1 & 33459 & 2156027 & 16922729 & 44372452 & 62471842 & \g59434367 & \g41336884 & \g21039624 & \g7487559 & \g1673055 & \g177147\\
\bottomrule
\end{tabular}
\caption{Fixed polyplets by height, $T(n,H)$. Shaded cells have $H > n/2$
and are given by the closed forms.}
\label{tab:tnh}
\end{table}

\paragraph{Non-fixed polyplets} Table~\ref{tab:onefree} and
Table~\ref{tab:bisym} give new counts for one-sided, free, bilateral and
asymmetric polyplets of size $n$, and for the free polyplets that are not
polyominoes (\oeis{A194596}), which are the free polyplets less the free
polyominoes (\oeis{A000105}). Few polyplets are polyominoes, since the
polyomino growth constant is at most $4.65$ \cite{klarnerRivest1973} against
the lower bound $6.22$ for polyplets below; and the one-sided count approaches
$a(n)/4$ and the free count $a(n)/8$ as $n$ grows, since a polyplet with a
nontrivial symmetry is determined by half of its cells and such polyplets are
exponentially rare (Section~\ref{sec:symmetry}). At $n = 32$ the one-sided
count exceeds $a(n)/4$ and the free count $a(n)/8$ by less than one part in
a million.

\begin{table}[htbp]
\centering
\begin{tabular}{r | r r}
\toprule
$n$         &                         one-sided &                            free\\
\midrule
\textbf{18} &               \num{5617792259411} &             \num{2808898025438}\\
\textbf{19} &              \num{37902303297525} &            \num{18951156321090}\\
\textbf{20} &             \num{256393397108358} &           \num{128196711128365}\\
\textbf{21} &            \num{1738521118535082} &           \num{869260590388450}\\
\textbf{22} &           \num{11813833328245848} &          \num{5906916748001325}\\
\textbf{23} &           \num{80437315244081648} &         \num{40218657830371643}\\
\textbf{24} &          \num{548666699140240976} &        \num{274333350132318510}\\
\textbf{25} &         \num{3748702832087014098} &       \num{1874351417444073722}\\
\textbf{26} &        \num{25651878467205504568} &      \num{12825939237386225482}\\
\textbf{27} &       \num{175781698028751634291} &      \num{87890849023826053192}\\
\textbf{28} &      \num{1206147307126536440324} &     \num{603073653588819503780}\\
\textbf{29} &      \num{8286282451482505589564} &    \num{4143141225805220756331}\\
\textbf{30} &     \num{56992306779773178547095} &   \num{28496153390059675197641}\\
\textbf{31} &    \num{392407904770584270738465} &  \num{196203952385726340537921}\\
\textbf{32} &   \num{2704550994366217058008641} & \num{1352275497184284183618433}\\
\textbf{33} &  \num{18657870495104684580672401}\\
\textbf{34} & \num{128829209605977867230343869}\\
\bottomrule
\end{tabular}
\caption{One-sided and free polyplet counts.}
\label{tab:onefree}
\end{table}

\begin{table}[htbp]
\centering\small
\begin{tabular}{r | r r r}
\toprule
$n$         &           bilateral &                      asymmetric &                 non-polyominoes\\
\midrule
\textbf{18} &       \num{3791465} &             \num{2808894233973} &             \num{2808705403386}\\
\textbf{19} &       \num{9344655} &            \num{18951146976435} &            \num{18950413696858}\\
\textbf{20} &      \num{25148372} &           \num{128196685979993} &           \num{128193840456415}\\
\textbf{21} &      \num{62241818} &           \num{869260528146632} &           \num{869249467327772}\\
\textbf{22} &     \num{167756802} &          \num{5906916580244523} &          \num{5906873556143637}\\
\textbf{23} &     \num{416661638} &         \num{40218657413710005} &         \num{40218489783363915}\\
\textbf{24} &    \num{1124396044} &        \num{274333349007922466} &        \num{274332695132618107}\\
\textbf{25} &    \num{2801133346} &       \num{1874351414642940376} &       \num{1874348860217028958}\\
\textbf{26} &    \num{7566946396} &      \num{12825939229819279086} &      \num{12825929238297403407}\\
\textbf{27} &   \num{18900472093} &      \num{87890849004925581099} &      \num{87890809870815114705}\\
\textbf{28} &   \num{51102567236} &     \num{603073653537716936544} &     \num{603073500077718909177}\\
\textbf{29} &  \num{127935923098} &    \num{4143141225677284833233} &    \num{4143140623183267694353}\\
\textbf{30} &  \num{346171848187} &   \num{28496153389713503349454} &   \num{28496151021712637626389}\\
\textbf{31} &  \num{868410337377} &  \num{196203952384857930200544} &  \num{196203943068019810549971}\\
\textbf{32} & \num{2351309228225} & \num{1352275497181932874390208} & \num{1352275460489267191905554}\\
\bottomrule
\end{tabular}
\caption{Counts of bilateral and asymmetric free polyplets, and the number of
non-polyomino free polyplets. Bilateral plus asymmetric equals free.}
\label{tab:bisym}
\end{table}

\paragraph{Holes} Polyplets with holes were tracked and counted during
transfer-matrix enumeration by carrying the Euler characteristic, through
$n = 18$. Redelmeier enumeration with a flood fill of each animal's
complement checked these counts for $n \le 14$. Counts are given in
Table~\ref{tab:holes}.

\begin{table}[H]
\centering\small
\begin{tabular}{r | r r r r}
\toprule
$n$ & $k=0$ & 1 & 2 & 3\\
\midrule
 1 & 1\\
 2 & 4\\
 3 & 20\\
 4 & 109 & 1\\
 5 & 622 & 16\\
 6 & \num{3664} & 166 & 2\\
 7 & \num{22094} & \num{1456} & 42\\
 8 & \num{135609} & \num{11788} & 538 & 6\\
 9 & \num{843941} & \num{91300} & \num{5596} & 144\\
10 & \num{5310754} & \num{688034} & \num{52268} & \num{2082}\\
11 & \num{33724862} & \num{5091820} & \num{457754} & \num{24144}\\
12 & \num{215793158} & \num{37209939} & \num{3841776} & \num{248526}\\
13 & \num{1389673091} & \num{269452496} & \num{31290820} & \num{2374520}\\
14 & \num{8998648488} & \num{1937956658} & \num{249306876} & \num{21556520}\\
15 & \num{58548155506} & \num{13865448048} & \num{1953275736} & \num{188568632}\\
16 & \num{382526638033} & \num{98796744112} & \num{15103433472} & \num{1603972074}\\
17 & \num{2508473632910} & \num{701664375696} & \num{115556223626} & \num{13348948048}\\
18 & \num{16503616943998} & \num{4970078092354} & \num{876484089730} & \num{109173877126}\\
\bottomrule
\end{tabular}\\[2ex]
\begin{tabular}{r | r r r r r r r}
\toprule
$n$ & $k=4$ & 5 & 6 & 7 & 8 & 9 & 10\\
\midrule
 9 & 1\\
10 & 42\\
11 & 820 & 8\\
12 & \num{11482} & 263 & 2\\
13 & \num{135235} & \num{4808} & 96\\
14 & \num{1437090} & \num{67020} & \num{2186} & 22\\
15 & \num{14261808} & \num{804084} & \num{35086} & 808 & 6\\
16 & \num{134769521} & \num{8779911} & \num{465138} & \num{16410} & 314 & 1\\
17 & \num{1227552608} & \num{89826208} & \num{5490242} & \num{251104} & \num{7847} & 104\\
18 & \num{10866411352} & \num{876032532} & \num{59950610} & \num{3272674} & \num{137298} & \num{3460} & 30\\
\bottomrule
\end{tabular}
\caption{Count of fixed polyplets of size $n$ with exactly $k$ holes.}
\label{tab:holes}
\end{table}

\paragraph{Growth rate} Placing one polyplet diagonally off the corner of
another's bounding box joins them into a polyplet from which both are
recoverable, so $a(m+n) \ge a(m)\,a(n)$, and by Fekete's lemma
%% JP:  I don't understand Fekete's lemma; I'd prefer that it come out of this document
$a(n)^{1/n}$ converges to $\lambda = \sup_n a(n)^{1/n}$, which is finite
because animals containing a given cell are bounded by rooted spanning trees
\cite{klarner1967}. Hence, from Table~\ref{tab:an} alone,
\[ \lambda \ge a(40)^{1/40} = 6.2208, \]
and $a(n) \le \lambda^n$ for every $n$. The ratios $a(n+1)/a(n)$ converge
to the same $\lambda$ \cite{madras1999}; the last is $a(40)/a(39) = 6.9352$
and they are still rising. A differential-approximant analysis of the first
$36$ terms estimates
\[ \lambda = 7.110(1), \qquad \theta = -1.000(1) \]
in $a(n) \sim A\,\lambda^n n^{\theta}$, the exponent matching the universal
value for two-dimensional lattice animals; the same code returns $4.06257$
for the polyomino constant against the accepted $4.0625696$. Rigorous bounds
$6.543 \le \lambda \le 9.3154$ are proved by exact certificates in a
companion paper \cite{L3}.

% ===========================================================================
\section{Methods}
\label{sec:methods}

\paragraph{Redelmeier enumeration} Redelmeier's algorithm
\cite{redelmeier1981} generates each fixed animal exactly once by growing
from a scan-minimal cell through an untried set. The implementation records
each animal's bounding box, so it produces the triangle $T(n,H)$ rather than just
$a(n)$, and with a flood fill of the complement, it counts holes. Its cost is
proportional to the number of animals, so it was used to compute $a(19)$, in
two independent campaigns on two machines, and to confirm the transfer
matrix's counts up to $n = 22$; the confirmation of the row $n = 22$ took
about $120$ hours on each of three machines, split into $24\,000$ disjoint
subtrees. Beyond $n = 22$ it cannot follow.

\paragraph{Transfer matrix method} Each height $H$ is an independent job. A
strip of height $H$ is swept column by column; a partial polyplet is
summarised by a signature recording, for each row of the most recent column,
%% JP:  American spellings please
whether the cell is occupied and if so which connected component of the
partial animal it belongs to, together with two flags for whether the top
and bottom rows of the strip have been reached. Two partial animals with the
same signature have the same completions, so counts are carried on
signatures. Each signature carries a vector of counts indexed by the number
of cells placed so far; combining two records adds the vectors,
%% and since addition is commutative and associative
%% JP:  wow really???!!
the result does not depend on how
the work was split across cores or machines. A signature is dropped when the
cells placed plus a lower bound on the cells still needed to reach both ends
of the strip and to join the components exceeds $n$. A signature whose next
column is empty, whose cells form one component and whose flags are both set
is one fixed polyplet of height exactly $H$. Because king adjacency lets
components cross, the non-crossing signature encodings used for polyominoes
\cite{jensen2001} do not apply.

Two transition kernels were used, and checked against each other on every
cell of the rows $n \le 29$. The first transfers a whole column per step,
resolving connectivity by union--find over the new column's cells and the old
%% JP:  union-find, surely?
column's components. The second advances one cell per step, as the polyomino
transfer matrices do, keeping the one already-overwritten north-west
neighbour a new cell may still attach to as a carried byte in the signature;
it is between $11\times$ and $55\times$ faster at the heights measured and
carried $a(30)$--$a(41)$. Counts are native $64$-bit integers through
$a(25)$ and $128$-bit from $a(26)$, with every addition guarded against overflow.

Several heights share one core pool, so that one height's merge phase hides
behind another's map phase; $a(40)$ ran as three phases on an $80$-core
machine, because its three tallest heights did not fit in memory together,
for $5.3$ million CPU-seconds and a peak of $363$~GB of disk.

\paragraph{Evaluating the tall heights} From $a(24)$ on, the heights with
$H > n/2$ and $k = n - H$ in the wired table were evaluated from $P_k$ rather
than swept, the table growing as levels were fixed: $k \le 7$ at $a(24)$ and
%% JP: "swept"/"sweep" is a terrible verb/noun.  you mean somethig specific by
%% it but I guarantee no-one will know what.  replace it with something less
%% fancy.
$k \le 18$ at $a(40)$, where heights up to $21$ were still swept. The sweep
of $a(40)$'s height $21$ alone took $36$ hours on $32$ cores; the tall
heights are the expensive ones because a height's cost follows its number of
signatures, not its count.

\paragraph{The colouring transfer matrix} The second source never decides
connectivity. For each $n$ it computes the sum over all $n$-cell subsets $S$
of a height-$H$ strip of $q^{c(S)}$, $c(S)$ the number of king-connected
components of $S$, by a frontier dynamic program over colour-coincidence
partitions in $\mathbb{Z}[q]/(q^2)$: an occupied cell with no earlier
occupied neighbour either adopts a colour already live or takes a fresh one
at weight $q - b$, and a cell whose earlier neighbours carry two different
colours contributes zero. The coefficient of $q$ is the number of connected
subsets. This is the Fortuin--Kasteleyn representation of the Potts partition
function specialised to site clusters \cite{fortuinKasteleyn1972}; the proof
in the engine's own terms is \repo{docs/proofs/cutcount-identity.md}. A
polyplet of height $h$ sits in a strip of height $H$ in $H - h + 1$
positions, so the strip count is $C_H(n) = \sum_{h \le H} (H-h+1)\,T(n,h)$
and $T(n,H) = C_H(n) - 2C_{H-1}(n) + C_{H-2}(n)$. The program runs one
prime at a time and reassembles exact counts by the Chinese remainder
theorem, with one prime always held out and predicted as a check. It shares
no code with either engine above.

%% JP:  Read to here
\paragraph{Symmetry counts}
\label{sec:symmetry}
By Burnside's lemma \cite{stanley2012}, with $\mathrm{Fix}_g(n)$ the number
of fixed polyplets of size $n$ invariant under $g$ in the symmetry group
$D_4$,
\begin{align*}
  \text{one-sided}(n) &= \tfrac14\bigl(a(n) + \mathrm{Fix}_{r_{180}}(n)
      + 2\,\mathrm{Fix}_{r_{90}}(n)\bigr), \\
  \text{free}(n) &= \tfrac18\bigl(a(n) + \mathrm{Fix}_{r_{180}}(n)
      + 2\,\mathrm{Fix}_{r_{90}}(n) + 2\,\mathrm{Fix}_{h}(n)
      + 2\,\mathrm{Fix}_{d}(n)\bigr),
\end{align*}
$h$ an axis reflection and $d$ a diagonal one, and the bilateral count is the
number of free classes containing a mirror-symmetric polyplet, obtained from
the same fixed-point counts by inclusion and exclusion over the reflection
subgroups. The fixed-point counts were computed by two programs, a direct
counter of symmetric animals and a transfer matrix restricted to them, which
agree where they overlap. A polyplet invariant under a nontrivial $g$ is
determined by its cells in a fundamental domain of $g$, so
$\mathrm{Fix}_g(n)$ is at most a polynomial in $n$ times $a(m)^2$ for
$m \le n/2 + 1$, while $a(n) \ge 6.2208^{\,n-39}$ and $a(m) \le \lambda^m$
with $\lambda \le 9.3154$; the ratio $\mathrm{Fix}_g(n)/a(n)$ therefore
decays exponentially, which is the claim about $a(n)/4$ and $a(n)/8$ in
Section~\ref{sec:results}.

\paragraph{Burnside congruences} Grouping the fixed polyplets by the exact
subgroup of $D_4$ that stabilises them, and writing $I(K)$ for the number
invariant under every element of a subgroup $K$, orbits of sizes $8$ and $4$
vanish modulo $4$ and the rest give
\[ a(n) \equiv I(C_4) + I(D_2^{\mathrm{ax}}) + I(D_2^{\mathrm{diag}})
   - 2\,I(D_4) \pmod 4, \]
with $C_4$ the rotations, $D_2^{\mathrm{ax}}$ the identity, both axis
reflections and the half turn, and $D_2^{\mathrm{diag}}$ likewise with the
diagonal reflections. Since $D_2^{\mathrm{ax}}$ is exactly the subgroup
preserving height, the same argument on the polyplets of one height gives
$T(n,H) \equiv I_H(D_2^{\mathrm{ax}}) \pmod 2$ for every cell. The invariant
counts are computed by an unrelated algorithm on the quotient domain in
minutes. The congruence holds for every $a(n)$, $n \le 40$, and the parity
for every one of the $820$ cells of the $n \le 40$ triangle.

\paragraph{Holes by the Euler characteristic} With king-adjacent foreground
and edge-adjacent background, the Euler number of a cell set is its number
of components minus its number of holes, and it is a sum over $2 \times 2$
windows of a local contribution \cite{gray1971}. Each window lies within one
pair of adjacent columns, so the Euler number accumulates over exactly the
column transitions the sweep already performs, as an additive coordinate
beside the cell count; a completed polyplet has $1 - E$ holes.

\paragraph{Formulas}
\label{sec:formulas}
A polyplet of height $H$ occupies every row of its bounding box, and since
king adjacency changes the row by at most one, a row holding a single cell is
a cut: every path from below it to above it passes through that cell.

$T(n,n)$: one cell per row, consecutive cells adjacent, so after fixing the
first cell's column each next cell is left, centre or right of the one
before it, $3^{\,n-1}$ in all.

$T(n,n-1)$: there are $n-1$ rows and $n-2$ junctions between rows where a
choice is made for the placement of the next cell or cells. A single cell has
exactly three options. To place $n$ cells in $n-1$ rows, one row must have
two cells, both connected to the rest: a joiner that is either a left--right
domino or a split, two cells separated by exactly one empty cell, since a
larger gap could not be bridged by the single cell of a neighbouring row. The
joiner may sit at any of the $n-3$ interior rows. For a domino joiner one
junction has $4$ possibilities rather than $3$ and the other also $4$, since
the domino may connect diagonally or orthogonally to the left or right of the
previous cell and the next cell may be placed likewise: $16(n-3)\,3^{\,n-4}$.
For a split joiner, when the split sits exactly left and right of the
previous cell there are five placements for the next cell, and otherwise
there are four placements for the split that each force one placement of the
next cell: $(4 \cdot 1 + 1 \cdot 5)(n-3)\,3^{\,n-4}$. Together $25(n-3)\,
3^{\,n-4}$ for an interior joiner. A joiner at the top or bottom has five
placements, giving $2 \cdot 5 \cdot 3^{\,n-3} = 30 \cdot 3^{\,n-4}$, and
\[ T(n,n-1) = 25(n-3)\,3^{\,n-4} + 30 \cdot 3^{\,n-4} = (25n - 45)\,3^{\,n-4}. \]
Each of the four weights $16$, $9$, $4$, $1$ was confirmed by enumerating
every fixed polyplet of size $n \le 8$ and classifying its doubled row, and
the formula holds on all $38$ banked cells $3 \le n \le 40$.

$T(n,n-2)$: either one row of three or two rows of two. A row of three has two
internal gaps, each $1$ or $2$ by the same argument; summing over gap types
and neighbour placements gives an interior weight of $49$ and an end weight
of $7$, so $(49n - 154)\,3^{\,n-5}$. Two doubled rows that are not adjacent
are two independent joiners of weight $25$, or $5$ at an end, and summing over
their placements gives $\tfrac12(625n^2 - 5375n + 11700)\,3^{\,n-7}$. Two
adjacent doubled rows form one compound joiner of weight $339$, or $66$ at an
end, giving $(1017n - 3897)\,3^{\,n-7}$. The three parts sum to
$\tfrac12(625n^2 - 2459n + 1134)\,3^{\,n-7}$; each part separately
matches the exhaustive count for $5 \le n \le 8$, and the sum holds on all
$36$ banked cells $5 \le n \le 40$ (\repo{docs/proofs/T-n-nm2-and-general.md}).

The general law (Theorem~\ref{thm:diag}) is the same bookkeeping carried to
all orders. Read a polyplet row by row; rows with one cell are walk rows and
maximal runs of rows with two or more cells are clusters, a cluster of
surplus $k$ having at most $k$ rows. A walk row is a cut, so a polyplet
decomposes uniquely into an optional bottom cluster, an alternating chain of
clusters and walk rows, and an optional top cluster. Counting each piece
relative to its walk cell, the generating function $F(y,z) = \sum T(n,H)\,
y^{\,n-H} z^{H}$ satisfies $F = E_b (1 - S)^{-1} E_t + P$ with $S = 3z +
\sigma$, where $\sigma$, $E_b$, $E_t$, $P$ are the cluster sums and every
term of $\sigma$ has positive $y$-order. Extracting the coefficient of $y^k$
involves finitely many clusters and at most $k$ factors of $\sigma$, so it is
a rational function of $z$ with denominator $(1 - 3z)^{k+1}$ plus a
polynomial of bounded degree, which is a polynomial of degree at most $k$ in
$H$ times $3^H$ once $H \ge k+1$, that is $n \ge 2k+1$. Resumming the chain
in one mode shows that each level $k$ adds exactly two new constants, $P_k(n)
= [y^k] \exp \sum_j (a_j + b_j n) y^j$, from which the leading coefficient
$25^k/k!$ follows from $P_1 = 25n - 45$. The full proofs are
\repo{docs/proofs/diagonal-law.md} and \repo{docs/proofs/grand-form.md}, and
the theorem is formalised in Lean~4.

% ===========================================================================
\section{Reproducibility}
\label{sec:reproduce}

No reader can re-run the computations, so this section says what stands
behind each number. A value computed twice by programs sharing no code is
confirmed. A value computed once and found consistent with proved structure
is checked, which catches gross errors and not a subtle systematic one.

\paragraph{Anchors} $a(1)$--$a(18)$ equal the eighteen terms the OEIS serves
for A006770, from both enumeration programs. Pointed at other classes, the
same machinery reproduces published integers it was never told: the fixed
polyominoes A001168 through their published range, Bacher's directed king
animals \cite{bacher2015}, and the polyomino growth constant from the
differential-approximant code (\repo{docs/external-anchors.md}).

\paragraph{Every run re-derives everything below it} Each production run
computes the whole triangle for its target size rather than extending a
stored one, and every earlier term is compared byte for byte. The $a(41)$
sweep reproduces the banked triangle at every one of the $610$ cells of
heights $1$--$20$ in rows $1$--$40$.

\paragraph{Second sources} Every cell of every row $n \le 22$ agrees between
Redelmeier enumeration and the transfer matrix. The colouring transfer matrix,
run to height $19$ at the size row $41$ needs, agrees with the production
engine on all $589$ cells with $n \le 40$ and $H \le 19$ and on the $19$
swept cells of row $41$, the held-out prime predicting every residue. A
third program, a strip transfer matrix with the same connectivity idea as the
production engine but none of its code, agrees on all $469$ cells with $H
\le 14$.

\paragraph{The closed forms} Each wired $P_k$ was fixed from two cells and
then met later real-swept cells it had not seen: $P_{15}$ at $T(33,18)$,
$P_{16}$ at $T(35,19)$, $P_{17}$ at $T(37,20)$, $P_{18}$ at $T(39,21)$.
Refitting every level from its two earliest cells above onset reproduces all
$171$ other real-swept cells above onset. Below the onset a cell obeys
$T(2k+1-j,\,k+1-j) = P_k(2k+1-j)\,3^{-(k+j)} + D_j(k)$ with a correction
$D_j(k)$ computed from the clusters of bounded surplus without reading the
triangle (\repo{results/undertow.md}); fixing every level $k \le 19$ from
cells of height at most $19$, all of which the colouring engine covers,
reproduces every cell of every row $n \le 39$ and re-sums each row to the
published term with no input from the production engine. That is the sense
in which $a(23)$--$a(39)$ are confirmed.

\paragraph{What stands behind each term}

\begin{table}[H]
\centering\small
\begin{tabular}{l R{0.72\textwidth}}
\toprule
terms & evidence \\
\midrule
$a(1)$--$a(18)$ & published; reproduced by both enumeration programs \\
$a(19)$--$a(22)$ & two algorithms sharing no code agree on every cell \\
$a(23)$--$a(39)$ & every cell enumerated by the colouring engine or given by a
  proved formula fixed from its cells; every row re-sums to the published
  term \\
$a(40)$ & heights $1$--$19$ two-source; heights $20$--$40$ from the formulas
  fixed on the colouring engine's cells; $T(40,20)$ and $T(40,21)$ enumerated
  once \\
$a(41)$ & heights $1$--$19$ two-source; height $20$ enumerated once and equal
  to the formula prediction; heights $21$--$41$ from the formulas, level $20$
  fixed below onset \\
\bottomrule
\end{tabular}
\caption{What stands behind each term.}
\label{tab:tiers}
\end{table}

Over the $820$ cells of the $n \le 40$ triangle, $628$ are enumerated by at
least one second program, $189$ are wired-formula cells above height $21$,
and three, $T(39,20)$, $T(40,20)$ and $T(40,21)$, have one enumeration and
no wired closed form. Those three are checked by the Burnside congruence and
parity and are reproduced by the formulas fixed on the colouring engine's
cells; none has been enumerated twice.

\paragraph{The forty-first term} $a(41)$ is a real sweep of heights
$1$--$20$ at size $41$, heights $1$--$19$ in $4.7$ hours on $40$ cores and
height $20$ in $9.6$ hours on $76$, plus the formulas for heights
$21$--$41$, with levels $k \le 19$ from the wired table and level $20$ fixed
from cells below its onset. The nineteen lowest swept cells agree with the
colouring engine, which has not been run at height $20$; that run is priced
at about a day and $100$~GB of memory. Before height $20$ was swept
the tower predicted $T(41,20)$ from level $21$, fixed three ways from
$T(38,17)$, $T(39,18)$ and $T(40,19)$ with the three fits agreeing; the sweep
returned exactly that value, $T(41,20) = \num{18004779862205054677763902712770}$, so the tower's top level has met an enumeration
it did not see, and $a(41)$ no longer uses it (\repo{make
gate-undertow-pairs}). A re-derivation built to differ wherever it could, with
its own solver and assembly and a different choice of held-out cells,
reproduces the term digit for digit (\repo{results/undertow-review-B.md}).
What remains: $T(41,20)$ has one enumeration and one formula prediction behind
it, not two enumerations, and heights $21$--$41$ are formula.

\paragraph{Recipes} The repository holds every banked number beside the run
record that produced it, and checks a reader can run in minutes. Each is
red-first, carrying a control that must fail.

\begin{center}\small
\begin{tabular}{l R{0.5\textwidth}}
\toprule
command & what it pins \\
\midrule
\repo{make gate-bfiles} & $a(1)$--$a(18)$ and the symmetry b-files against the OEIS \\
\repo{make gate-g2} & the Redelmeier engine against an oracle and fixtures \\
\repo{make gate-kink-oracle} & the cell-at-a-time kernel, every height real, against $a(1)$--$a(18)$ \\
\repo{make gate-cutcount-assembly} & the colouring engine's rows assemble to the triangle; held-out primes \\
\repo{make ns-gate-diag-pins} & every wired $P_k$ against every real-swept cell above onset \\
\repo{make gate-undertow-pairs} & the $a(41)$ tower, level $21$ three ways, and the swept $T(41,20)$ against its prediction \\
\repo{make gate-subgroup} & the Burnside congruence \\
\repo{make gate-provenance} & the per-cell source table and its counts \\
\repo{python3 paper/verify\_technical\_report.py} & every number in the tables against the banked results \\
\bottomrule
\end{tabular}
\end{center}

Each term's run record is the \repo{PROVENANCE.md} in its
\repo{results/ns\_a\textit{NN}/} directory, $19 \le NN \le 40$, and in
\repo{results/a41/}; the triangle is
\repo{results/triangle.txt}; the symmetry fixed-point counts
\repo{results/sym\_counts.txt}; the hole counts \repo{results/holes\_n18.txt}
and \repo{results/holes\_n14.txt}; the Lean development \repo{polyplets/},
whose \repo{PROOF-STATUS.md} says what it covers and what it does not. The
plain-language account of how well each value is supported, kept current, is
\repo{results/confidence.md}.

\bibliographystyle{plain}
\bibliography{technical-report}

\end{document}
